\begin{problem}{}{}{UCB EECS127 FA23 - Disc 8}
In this question we will see how combinatorial optimization problems can sometimes be solved
via related convex optimization problems. Consider the problem of deciding which subset of items
to sell from among a collection of $n$ items. Selling the item $i$ results in revenue $s_i > 0$
and a transaction cost $c_i > 0$. There is also a fixed overall cost, which we normalize to $1$,
which is incurred irrespective of whether any items are sold. We are interested in maximizing the \emph{margin},
i.e. the ratio of the total revenue to the total cost (sum of the total transaction cost and the fixed overall cost).
Note that this is a combinatorial optimization problem, because the set of choices is a discrete set
(here it is the set of all subsets of $\{1, \dots, n\}$). In this question we will see that it is possible to
pose this combinatorial optimization problem as a convex optimization problem.

\begin{enumerate}
    \item[(a)] Show that the original combinatorial optimization problem can be formulated as
          \[
              \max_{\vec{x} \in \{0,1\}^n} f(\vec{x})
          \]
          where $f(\vec{x}) := \frac{\vec{s}^\top \vec{x}}{1 + \vec{c}^\top \vec{x}}$, Here $\vec{s} \in \mathbb{R}^n$ is the column
          vector of revenues and $\vec{c} \in \mathbb{R}^n$ is the column vector of transaction costs associated
          to the individual items.

    \item[(b)] Show that the combinatorial optimization problem admits a convex optimization problem formulation
          in the sense that the convex optimization problem
          \begin{align*}
              \min_t \quad             & t                                                   \\
              \text{subject to:} \quad & t \ge \vec{\mathbf{1}}^\top (\vec{s} - t\vec{c})_+, \\
                                       & t \ge 0,
          \end{align*}
          has the same value as that of the original combinatorial optimization problem. Here $\vec{z}_+$,
          for a vector $\vec{z} \in \mathbb{R}^n$, denotes the vector with components $\max(0, z_i)$, $i = 1, \dots, n$.
          Also $\vec{\mathbf{1}}$ denotes the all-ones column vector of length $n$.

          \textit{Note:} You should also justify why the problem formulated in this part of the question is
          a convex optimization problem.

          \textit{Hint:} For given $t \ge 0$, express the condition ``$f(\vec{x}) \le t$ for every
          $\vec{x} \in \{0,1\}^n$''
          in simple terms. Note that this is related to the idea of introducing slack variables,
          as in Sec. 8.3.4.4 of the
          textbook of Calafiore and El Ghaoui.

    \item[(c)] Show that the inequality constraint $t \ge \vec{\mathbf{1}}^\top (\vec{s} - t\vec{c})_+$ is active at the
          optimum in the convex optimization problem, namely, it will have to be satisfied with equality.

    \item[(d)] How can you recover an optimal solution $\vec{x}^*$ for the original combinatorial optimization problem
          from an optimal value $t^*$ for the above convex optimization problem?
\end{enumerate}
\end{problem}

\begin{solution}
    \begin{itemize}

        \item[(a)] Maximizing the margin:
              \begin{equation}
                  \max_{S \in \mathcal{P}(\{1,\dots,n\})}
                  \frac{\text{total revenue of selling items in subset } S}%
                  {1 + \text{total transaction cost of selling items in subset } S} 
              \end{equation}

              Represents the subset $S$ as a binary mask in $\mathbb{R}^n$, we rewrite (1):
              \begin{align*}
                  \max_{\vec{x} \in \{0,1\}^n} f(\vec{x}), \quad
                  f(\vec{x}) & = \frac{\vec{s}^\top \vec{x}}{1 + \vec{c}^\top \vec{x}}
              \end{align*}

              Here $\vec{s} \in \mathbb{R}^n$ is the column vector of revenues
              and $\vec{c} \in \mathbb{R}^n$ is the column vector of transaction costs.

        \item[(b)] The problem $\max_{\vec{x} \in \{0,1\}^n} f(\vec{x})$ is equivalent
              to the problem:
              \begin{align*}
                  \min_t \quad             & t                                                         \\
                  \text{subject to } \quad & t \ge f(\vec{x}) \text{ for every } \vec{x} \in \{0,1\}^n
              \end{align*}

              Rewrite $t \ge f(\vec{x})$ for every $\vec{x} \in \{0,1\}^n$:
              \begin{align*}
                  t \ge \frac{\vec{s}^\top \vec{x}}{1 + \vec{c}^\top \vec{x}},
                  \; \forall \vec{x} \in \{0,1\}^n
                   & \Leftrightarrow t + \vec{x}^\top(t\vec{c}) \ge \vec{s}^\top \vec{x},
                  \; \forall \vec{x} \in \{0,1\}^n                                                       \\
                   & \Leftrightarrow t \ge \vec{x}^\top(\vec{s} - t\vec{c}),
                  \; \forall \vec{x} \in \{0,1\}^n                                                       \\
                   & \Leftrightarrow t \ge \max_{\vec{x} \in \{0,1\}^n} \vec{x}^\top(\vec{s} - t\vec{c}) \\
                   & \Leftrightarrow t \ge \mathbf{1}^\top(\vec{s} - t\vec{c})_+
              \end{align*}

              Therefore,
              \begin{equation*}
                  \max_{\vec{x} \in \{0,1\}^n} f(\vec{x}) =
                  \begin{aligned}[t]
                       & \min t                                                   \\
                       & \text{s.t: } t \ge \mathbf{1}^\top(\vec{s} - t\vec{c})_+
                  \end{aligned}
              \end{equation*}

              The reformulated problem is a convex optimization problem since:
              \begin{itemize}
                  \item The objective function $f_0(t) = t$ is convex in $t$
                  \item The constraint $f_1(t) = -t + \max_{\vec{x} \in \{0,1\}^n}
                            \vec{x}^\top(\vec{s} - t\vec{c}) \le 0$ has $f_1$ be convex in $t$ since:
              \end{itemize}
              \begin{align*}
                  f_1(t) & = -t + \max_{\vec{x} \in \{0,1\}^n} \vec{x}^\top(\vec{s} - t\vec{c}) \\
                         & = -t + \sum_{i=1}^n \max(0, s_i - tc_i)
              \end{align*}

              Apply Nonnegative weighted sum and pointwise maximization rule
              $\Rightarrow f_1$ convex.

        \item[(c)] Let $t^*$ be the solution to the optimization problem.
              Suppose the inequality constraint is not active at the optimum.

              We have:
              \begin{align*}
                                    & t^* > \mathbf{1}^\top(\vec{s} - t^*\vec{c})_+ \\
                  \Rightarrow \quad & t^* > \max_{\vec{x} \in \{0,1\}^n}
                  \vec{x}^\top(\vec{s} - t^*\vec{c})                                \\
                  \Rightarrow \quad & t^* > \vec{x}^\top(\vec{s} - t^*\vec{c}),
                  \quad \forall \vec{x} \in \{0,1\}^n                               \\
                  \Rightarrow \quad & t^* > f(\vec{x}),
                  \quad \forall \vec{x} \in \{0,1\}^n                               \\
                  \Rightarrow \quad & t^* > \max_{\vec{x} \in \{0,1\}^n} f(\vec{x})
              \end{align*}

              Which contradicts the fact that $t^* = \max_{\vec{x} \in \{0,1\}^n} f(\vec{x})$.

        \item[(d)] Constraint is active at optimum:
              \begin{align*}
                  \Rightarrow \quad & t^* = \max_{\vec{x} \in \{0,1\}^n}
                  \vec{x}^\top(\vec{s} - t^*\vec{c})                                             \\
                  \Rightarrow \quad & \exists \vec{x}^* \in \{0,1\}^n :
                  t^* = \vec{x}^{*\top}(\vec{s} - t^*\vec{c})                                    \\
                  \Rightarrow \quad & \vec{x}^* = \operatorname*{argmax}_{\vec{x} \in \{0,1\}^n}
                  f(\vec{x})
              \end{align*}

              Since $\vec{x}^* = \operatorname*{argmax}_{\vec{x} \in \{0,1\}^n}
                  \vec{x}^\top(\vec{s} - t^*\vec{c})$, construct $\vec{x}$ such that:
              \begin{equation*}
                  x_i = \begin{cases}
                      1 & \text{if } s_i > t^*c_i   \\
                      0 & \text{if } s_i \le t^*c_i
                  \end{cases}
              \end{equation*}

    \end{itemize}
\end{solution}